\chapter{Zenoh}
\label{chap:zenoh}

\section{Cos'è Zenoh}
\label{sec:cose-zenoh}

Descrizione nel dettaglio di cosa è Zenoh
\begin{itemize}
	\item Modello di comunicazione peer-to-peer: scouting automatico dei nodi,
	differenza tra client, peer e router.
	\item Possibilità di usare sia un modello pub/sub che query/reply, quali sono
	le differenze e come possono essere usati entrambi in uno stesso programma.
\end{itemize}

\section{Zenoh Pico}
\label{sec:zenoh-pico}

Introduzione al problema dovuto all'impossibilità di utilizzare Zenoh su sistemi
embedded e come zenoh pico risolva questo problema. Piccola menzione anche di
zenoh-nostd, e di come però sia ancora in fase di sviluppo.
\begin{itemize}
	\item Descrizione di limitazioni e differenze di zenoh pico rispetto a zenoh.
	\item Descrizione dello \textquote{ownership model} imitato tramite funzioni C per
	mantenere la libreria il più simile possibile all'originale e per evitare
	più facilmente problemi dovuti all'uso improprio della memoria.
\end{itemize}
